% v10: independently audited modular consequences of the local equations
% and puncture clutching.  The modular results use no basin positivity or
% bounded search.  The last geometric corollary also uses qh9:isolated-edge.

\section{素数 support に対する五次の合同障害}
\label{ps10:section}

全平面五次の五つの cores の標識集合を $V=\{0,1,2,3,4\}$ とする．
共有 branch $ij$ の分母を $q_{ij}=q_{ji}$，core $i$ の branch $ij$ における
puncture $ijk$ の numerator を $p^i_{jk}$ と書く．
以下では $q_{ij}$ が正の奇整数で，$p^i_{jk}\in\ZZ$ が，相異なる $i,j,k$ に対し
\begin{align}
 \sum_{k\in V\setminus\{i,j\}}p^i_{jk}&=-q_{ij},
       \label{ps10:row}\\
 p^i_{jk}p^i_{kj}&=q_{ij}q_{ik}-1,
       \label{ps10:det}\\
 q_{ij}&\mid p^i_{jk}+p^j_{ik}
       \label{ps10:clutch}
\end{align}
を満たすと仮定する．最初の二式は局所の三 puncture period relation と
crossing unimodularity，第三式は共有 puncture cycles の整数 clutching から来る．
以下の合同障害はこの三条件だけを用い，basin 面積の正値性は用いない．
後の系~\ref{ps10:all-high-divisible-three} では，さらに既存の幾何的接着障害を用いる．

奇素数 $\ell$ ごとに，頂点集合 $V$ 上の単純 graph
\[
 E(G_\ell)=\bigl\{\{i,j\}:\ell\mid q_{ij}\bigr\}
\]
を定める．部分 graph を含むという表現は induced であることを要求しない．
以下の合同計算はすべて $\mathbf F_\ell$ 内で行う．
$ij\in E(G_\ell)$ なら，\eqref{ps10:det} により，すべての $k\ne i,j$ で
\begin{equation}\label{ps10:unit}
 p^i_{jk}\ne0,
 \qquad p^i_{kj}=-(p^i_{jk})^{-1}
\end{equation}
となることを繰り返し用いる．

\begin{lemma}[三角形の符号]\label{ps10:triangle}
$i,j,k$ が $G_\ell$ の三角形をなすなら，ある
$\varepsilon\in\{1,-1\}\subset\mathbf F_\ell$ に対し
\[
 p^i_{jk}=p^j_{ki}=p^k_{ij}=\varepsilon,
 \qquad
 p^i_{kj}=p^j_{ik}=p^k_{ji}=-\varepsilon
\]
である．特に，この三角形に属する各 numerator の二乗は一である．
\end{lemma}

\begin{proof}
$x=p^i_{jk}\ne0$ と置く．core $i$ の determinant と，辺 $ij,ik$ の
clutching を順に用いると
\[
 p^i_{kj}=-x^{-1},\qquad p^j_{ik}=-x,\qquad p^k_{ij}=x^{-1}
\]
を得る．core $j$ では $p^j_{ki}=x^{-1}$ であり，辺 $jk$ の clutching は
$p^k_{ji}=-x^{-1}$ を与える．core $k$ の determinant より
$-x^{-2}=-1$，すなわち $x^2=1$ となる．
$\ell$ は奇素数なので $x=\varepsilon=\pm1$ であり，表示した六式を得る．
\end{proof}

\begin{theorem}[素数 support の禁止配置]\label{ps10:forbidden}
\eqref{ps10:row}--\eqref{ps10:clutch} の下で，次が成り立つ．
\begin{enumerate}
\item 任意の奇素数 $\ell$ に対し，$G_\ell$ は $K_4$ を含まない．
\item $G_3$ は matching，すなわち最大次数が一以下の graph である．
\item 任意の奇素数 $\ell$ に対し，$G_\ell$ は book $B_3$ を含まない．
  ここで $B_3$ は一つの辺 $ij$ と，その両端に接続する三つの相異なる
  vertices $k,l,m$ からなる七辺の graph である．
\item $G_\ell$ が $K_1\vee P_4$ を含むなら，$\ell=5$ である．
  ここで $K_1\vee P_4$ は，path $a-b-c-d$ の四 vertices すべてに
  一つの vertex $o$ を接続した七辺の graph を表す．
\end{enumerate}
\end{theorem}

\begin{proof}
まず $K_4$ の vertices を $i,j,k,l$，残る vertex を $m$ とする．
core $i$ で，row $j$ の entries $p^i_{jk},p^i_{jl}$ は
補題~\ref{ps10:triangle} により $\pm1$ である．
これらが逆符号なら，row sum より $p^i_{jm}=0$ となり，
\eqref{ps10:unit} に反する．従って二つは同じ符号 $s_j$ を持つ．
同様に rows $k,l$ の二つの内部 entries も，各々同じ符号 $s_k,s_l$ を持つ．
各 pair の determinant は
\[
 s_js_k=s_js_l=s_ks_l=-1
\]
を要求するが，三つの $\pm1$ が互いに逆符号となることはできない．
これが第一項である．

第二項では $\ell=3$ とし，ある core に接する二本以上の edges が
$G_3$ に属すると仮定する．この core の四つの局所 branch labels を
$0,1,2,3$ と書き直し，$q_0=q_1=0$ in $\mathbf F_3$ としてよい．
rows $0,1$ の三つの entries はそれぞれ全て非零であり，row sum は零である．
$\mathbf F_3$ で三つの $\pm1$ の和が零であるためには，全てが同じ符号でなければならない．
従ってある $\sigma,\tau\in\{1,-1\}$ に対し
\[
 p_{01}=p_{02}=p_{03}=\sigma,
 \qquad p_{10}=p_{12}=p_{13}=\tau.
\]
pair $01$ の determinant は $\sigma\tau=-1$ を与えるので，$\tau=-\sigma$ である．
残る $j=2,3$ について，pairs $0j,1j$ の determinants より
\[
 p_{j0}=-\sigma,\qquad p_{j1}=\sigma.
\]
これらは各 row で打ち消し合うため，$p_{23}=-q_2$，$p_{32}=-q_3$ となる．
しかし pair $23$ の determinant は
\[
 q_2q_3=p_{23}p_{32}=q_2q_3-1
\]
を要求し，矛盾する．$q_2,q_3$ の剰余には追加条件を課していない．
従って一つの core に三で割れる分母は高々一つであり，$G_3$ は matching である．
この項では clutching 条件も面積正値性も使っていない．

第三項の $B_3$ があれば，core $i$ の row $j$ は三つの三角形 $ijk,ijl,ijm$ により
全 entries が $\pm1$ となる．その和が零であることから $\ell=3$ が必要である．
しかし $i$ は $G_3$ で次数四を持ち，第二項に反する．

最後に $K_1\vee P_4$ を仮定し，core $o$ の添字 $o$ を省略する．
三角形の補題から
\[
 p_{ab}=s,\qquad p_{bc}=t,\qquad p_{cd}=u,
 \qquad s,t,u\in\{1,-1\}
\]
と書ける．この core の四分母は全て零であり，すべての非対角 entries が非零である．
row $b$ より $p_{bd}=s-t\ne0$ なので，$t=-s$，$p_{bd}=2s$ となる．
row $c$ より $p_{ca}=t-u=-s-u\ne0$ なので，$u=s$，$p_{ca}=-2s$ となる．
determinant と row $a$ を用いると
\[
 p_{ac}=\frac{s}{2},\qquad p_{ad}=-\frac{3s}{2}.
\]
$p_{ad}\ne0$ より $\ell\ne3$ である．従って row $d$ は
\[
 p_{da}+p_{db}+p_{dc}
   =\frac{2s}{3}-\frac{s}{2}-s
   =-\frac{5s}{6}=0
\]
となり，$\ell=5$ が従う．
\end{proof}

\begin{proposition}[$K_{2,2,1}$ と残る二分母]\label{ps10:k221}
$G_\ell$ が parts $\{a,b\},\{c,d\},\{o\}$ の $K_{2,2,1}$ を含むなら，
$\ell=5$ であり，$G_5$ はこの $K_{2,2,1}$ に等しい．
さらに，ある $s\in\{1,-1\}\subset\mathbf F_5$ に対し
\begin{equation}\label{ps10:k221-residues}
 q_{ab}\equiv-s,\qquad q_{cd}\equiv s\pmod5.
\end{equation}
特に $q_{ab}+q_{cd}\equiv0$ および $q_{ab}^2\equiv q_{cd}^2\equiv1\pmod5$ である．
\end{proposition}

\begin{proof}
$K_{2,2,1}$ は $K_1\vee P_4$ を含むので，定理~\ref{ps10:forbidden} より
$\ell=5$ である．辺 $ab$ または $cd$ を加えると $K_4$ が生じるため，
$q_{ab},q_{cd}$ は五で割れず，support graph はちょうど $K_{2,2,1}$ となる．

まず core $o$ を考える．各 row の反対 part への二 entries は $\pm1$ であり，
残る一 entry が非零であるため，その二つは同符号である．
cross-part determinants から，part $\{a,b\}$ の二 rows は同じ符号 $s$，
part $\{c,d\}$ の二 rows は逆符号 $-s$ を持つ．
row sums より
\[
 p^o_{ab}=p^o_{ba}=-2s,
 \qquad p^o_{ac}=p^o_{ad}=p^o_{bc}=p^o_{bd}=s,
 \qquad p^o_{cd}=p^o_{dc}=2s.
\]

次に core $a$ を考える．辺 $ao$ の clutching と local determinants より
\[
 p^a_{ob}=2s,\qquad p^a_{oc}=p^a_{od}=-s,
 \qquad p^a_{bo}=2s,\qquad p^a_{co}=p^a_{do}=s
\]
となる．ここでは $-1/2=2$ in $\mathbf F_5$ を用いた．
$p^a_{cd}=s x$ と書くと，$x\ne0$ であり，rows $c,d$ と determinants は
\[
 p^a_{cb}=-s(1+x),\qquad
 p^a_{db}=s(-1+x^{-1}),\qquad
 p^a_{bc}=\frac{s}{1+x},\qquad
 p^a_{bd}=\frac{s}{1-x^{-1}}
\]
を与える．最初の二 entries が非零なので $x\ne\pm1$，従って $x=\pm2$ である．
row $b$ の和は
\[
 -q_{ab}
 =s\left(2+\frac1{1+x}+\frac1{1-x^{-1}}\right)
 =\begin{cases}s&x=2,\\0&x=-2.\end{cases}
\]
$q_{ab}\not\equiv0$ なので $x=2$ であり，$q_{ab}\equiv-s$ となる．
二つの parts を交換すると $s$ は $-s$ に変わるので，同じ議論から
$q_{cd}\equiv s$ を得る．
\end{proof}

\begin{corollary}[予備的な edge 数評価と共通因子]\label{ps10:counts-gcd}
\eqref{ps10:row}--\eqref{ps10:clutch} の下で
\[
 |E(G_\ell)|\le
 \begin{cases}
 2&\ell=3,\\
 8&\ell=5,\\
 7&\ell\ne3,5
 \end{cases}
\]
である．また，任意の四元部分集合 $S\subset V$ に対し
\begin{equation}\label{ps10:gcd}
 \gcd\{q_{ij}:i,j\in S,\ i<j\}=1.
\end{equation}
従って，全十本の共有 branches が同じ分母 $q>1$ を持つ配置は存在しない．
\end{corollary}

\begin{proof}
$G_3$ は五 vertices 上の matching なので，edge 数は二以下である．
五 vertices 上の $K_4$-free graph は八 edges 以下である．
実際，九 edges 以上なら欠ける edge は高々一本なので $K_4$ を含む．
八 edges の場合，欠ける二 edges が共通 vertex を持つと，残る四 vertices が
$K_4$ をなす．従って欠ける二 edges は disjoint であり，元の graph は
$K_{2,2,1}$ に等しい．この場合は命題~\ref{ps10:k221} により $\ell=5$ である．
これが三つの edge bounds を与える．

\eqref{ps10:gcd} の左辺が一より大きければ，その奇素因数 $\ell$ に対し
$S$ が $G_\ell$ 内の $K_4$ をなす．これは定理~\ref{ps10:forbidden} に反する．
全分母が同じ $q$ なら，任意の四元部分集合の gcd も $q$ なので最後の結論を得る．
\end{proof}

\begin{corollary}[全高分母が三で割れる場合]\label{ps10:all-high-divisible-three}
本稿の全平面五次の SNC core models，標準 $K_4$ cones，既存の branch graph を用い，
追加の branch singularities を入れない大域作用適合 atlas を考える．
すべての共有 branch に対して
\[
 q_{ij}>1\quad\Longrightarrow\quad 3\mid q_{ij}
\]
が成り立つなら，すべての $q_{ij}$ は一である．
同値に，高分母を持つ大域 atlas には，三で割れない高分母も必要である．
\end{corollary}

\begin{proof}
高分母 branches だけからなる support graph を $H$ とする．
仮定の下では $H=G_3$ であり，定理~\ref{ps10:forbidden} により $H$ は matching である．
$H$ が非空なら，その各 edge は単独で一つの連結成分をなす．
これは系~\ref{qh9:isolated-edge} に反する．
具体的には，その両端の core は分母 pattern $(q,1,1,1)$ を持ち，
実際の atlas の正の basin areas により，命題~\ref{qh9:isolated-local} が両端の
normal-action 分岐係数をともに負とする．
同じ neck の対向 normal circles から必要な係数の相殺は成立し得ない．
従って $H$ は空である．
\end{proof}

\begin{remark}[七辺の場合の残る graph]\label{ps10:seven-edge-shape}
$\ell\ne3,5$ で $|E(G_\ell)|=7$ なら，complement は $P_3\sqcup K_2$ に限られる．
同値に，$G_\ell$ は $K_{2,3}$ の三点側に一本の edge を加えた graph である．
実際，五 vertices 上の三辺 graph は，isolated vertices を省いて
$K_3,K_{1,3},P_4,P_3\sqcup K_2$ のいずれかである．
前の三つの complements は，順に $B_3$，$K_4$ を含む graph，$K_1\vee P_4$ なので排除される．
この段階の必要条件は，残る七辺 graph の実現可能性を主張しない．
次節の定理~\ref{ps11:dense} では，五つの cores の式を用いて七辺ちょうどの場合を全て排除する．
\end{remark}

以上は素数ごとの普遍的な合同障害であり，分母の大きさを制限した探索の結論ではない．
一方，異なる素因数を持つ分母の配置や，許された support graphs の整数 lifting，
basin areas と有限 neck の同時接着は，これらの必要条件だけでは決まらない．
特に，系~\ref{ps10:all-high-divisible-three} の追加仮定を外した任意の分母配置について，
全平面五次のすべての分母が一であることまでは結論しない．
